\documentclass{article}
\usepackage{amsmath,amssymb,amsthm}
\usepackage{algorithm,algorithmic}
\usepackage{hyperref}
\usepackage{enumitem}
\usepackage{geometry}
\geometry{margin=1in}
\newtheorem{theorem}{Theorem}
\newtheorem{lemma}{Lemma}
\newcommand{\E}{\mathbb{E}}
\newcommand{\R}{\mathbb{R}}
\newcommand{\norm}[1]{\left\|#1\right\|}
\newcommand{\ip}[2]{\left\langle #1,#2\right\rangle}
\newcommand{\block}[1]{\mathcal{B}_{#1}}
\newcommand{\U}[1]{U_{#1}}

\begin{document}

\section*{Random Block Coordinate Descent (RBCD)}

\section*{Mathematical Formulation}
Let $f:\R^{n}\to\R$ be a convex differentiable function.  
The coordinates are partitioned into $m$ disjoint blocks  
\[
\{1,\dots ,n\}=\bigcup_{i=1}^{m}\block{i},\qquad 
\block{i}\cap\block{j}=\emptyset\;(i\neq j),
\]
and let $U_i\in\R^{n\times n_i}$ be the column‑selection matrix that extracts block $i$ from a vector (so $U_i^{\top}U_i=I_{n_i}$ and $U_iU_i^{\top}$ projects onto the coordinates of block $i$).

For each block $i$ assume a \emph{block‑coordinate smoothness} constant $L_i>0$ such that for all $x\in\R^{n}$ and any $d\in\R^{n_i}$,
\begin{equation}
\label{eq:smooth}
f\bigl(x+U_i d\bigr)\le 
f(x)+\ip{\nabla_i f(x)}{d}+\frac{L_i}{2}\norm{d}^{2},
\end{equation}
where $\nabla_i f(x):=U_i^{\top}\nabla f(x)$ denotes the partial gradient w.r.t. block $i$.

Given a stepsize $\eta>0$, the Random Block Coordinate Descent iteration is:
\begin{enumerate}[label=(\roman*)]
\item Sample a block index $i_t\in\{1,\dots ,m\}$ uniformly at random,
      i.e. $\Pr(i_t=i)=1/m$.
\item Perform a gradient step on the selected block:
\begin{equation}
\label{eq:update}
w_{t+1}=w_t-\eta\,U_{i_t}\nabla_{i_t} f(w_t)
      =w_t-\eta\,U_{i_t}U_{i_t}^{\top}\nabla f(w_t).
\end{equation}
\end{enumerate}
All other coordinates remain unchanged.

\section*{Pseudocode}
\begin{algorithm}[H]
\caption{Random Block Coordinate Descent}
\label{alg:rbcd}
\begin{algorithmic}[1]
\REQUIRE Initial point $w_0\in\R^{n}$, stepsize $\eta>0$, maximum iterations $T$.
\FOR{$t=0$ \TO $T-1$}
    \STATE Sample $i_t\sim\text{Uniform}\{1,\dots ,m\}$.
    \STATE Compute the partial gradient $\displaystyle g_{i_t}= \nabla_{i_t}f(w_t)=U_{i_t}^{\top}\nabla f(w_t)$.
    \STATE Update the selected block:
    \[
    w_{t+1}= w_t-\eta\,U_{i_t} g_{i_t}.
    \]
\ENDFOR
\RETURN $w_T$ (or the average $\bar w_T=\frac1T\sum_{t=0}^{T-1}w_t$).
\end{algorithmic}
\end{algorithm}

\section*{Dry Run Example}
Consider the scalar problem $f(w)=(w-3)^2$, i.e. $n=1$, $m=1$ (a single block).  
Take $w_0=0$, stepsize $\eta=0.1$, and run three iterations.

\begin{center}
\begin{tabular}{c|c|c|c}
$t$ & $g_t=\nabla f(w_t)$ & $w_{t+1}=w_t-\eta g_t$ & $f(w_{t+1})$\\\hline
0 & $2(w_0-3)=-6$ & $0-0.1(-6)=0.6$ & $(0.6-3)^2=5.76$\\
1 & $2(0.6-3)=-4.8$ & $0.6-0.1(-4.8)=1.08$ & $(1.08-3)^2=3.6864$\\
2 & $2(1.08-3)=-3.84$ & $1.08-0.1(-3.84)=1.464$ & $(1.464-3)^2=2.3665$
\end{tabular}
\end{center}
The objective value decreases at each step, illustrating the algorithm.

\newpage

\section*{Assumptions}
\begin{enumerate}[label=(A\arabic*)]
\item \textbf{Convexity.} $f$ is convex:
\[
f(y)\ge f(x)+\ip{\nabla f(x)}{y-x},\qquad\forall x,y\in\R^{n}.
\tag{A1}
\]
\item \textbf{Block‑coordinate smoothness.} For each block $i$, \eqref{eq:smooth} holds with constant $L_i>0$.
\tag{A2}
\item \textbf{Bounded domain.} There exists $R>0$ such that for the (unknown) optimal point $w^{\ast}$,
\[
\norm{w_0-w^{\ast}}\le R .
\tag{A3}
\]
\item \textbf{Stepsize.} The stepsize satisfies
\[
0<\eta\le\frac{1}{L_{\max}},\qquad 
L_{\max}:=\max_{i}L_i .
\tag{A4}
\]
\end{enumerate}

\section*{Main Convergence Theorem}
\begin{theorem}
\label{thm:conv}
Under assumptions (A1)–(A4), let $\{w_t\}_{t\ge0}$ be the iterates generated by RBCD.  
Define the average iterate $\bar w_T:=\frac{1}{T}\sum_{t=0}^{T-1}w_t$.  
Then for any optimal solution $w^{\ast}$,
\[
\E\!\left[f(\bar w_T)-f(w^{\ast})\right]
\;\le\;
\frac{R^{2}}{2\eta T}.
\tag{T1}
\]
Consequently, the expected suboptimality decays as $O(1/T)$.
\end{theorem}

\section*{Theoretical Properties}
\begin{itemize}
\item \textbf{Stability.} The bound (T1) holds for any stepsize $\eta\le 1/L_{\max}$, guaranteeing that the algorithm does not diverge.
\item \textbf{Hyperparameter sensitivity.} A larger $\eta$ (up to $1/L_{\max}$) yields a smaller constant $R^{2}/(2\eta)$, i.e. faster convergence in practice; however, exceeding $1/L_{\max}$ invalidates (A4) and may break the descent property.
\item \textbf{Comparison to full‑gradient descent.} Full gradient descent with stepsize $\eta\le 1/L_{\text{full}}$ (where $L_{\text{full}}$ is the Lipschitz constant of $\nabla f$) enjoys the same $O(1/T)$ rate but each iteration costs $O(n)$. RBCD costs only $O(n_i)$ per iteration, offering computational savings for large $n$ when blocks are small.
\end{itemize}

\section*{Discussion}
The theorem shows that random block updates retain the optimal $O(1/T)$ rate for convex smooth problems, matching the classic analysis of stochastic gradient descent but with a deterministic per‑iteration computation on a subset of coordinates. Extensions to non‑uniform sampling, accelerated schemes, or strongly convex functions follow analogously and improve the constants or rates.

\newpage

\section*{Preliminaries (Definitions and Lemmas)}
\begin{lemma}[Expected squared distance decrease]
\label{lem:dist}
For any $w\in\R^{n}$ and any optimal $w^{\ast}$,
\[
\E_{i\sim\text{Unif}[m]}\bigl[\norm{w-\eta U_i\nabla_i f(w)-w^{\ast}}^{2}\bigr]
\le
\norm{w-w^{\ast}}^{2}
-2\eta\bigl(f(w)-f(w^{\ast})\bigr)
+\eta^{2}L_{\max}\norm{\nabla f(w)}^{2}.
\tag{L1}
\]
\end{lemma}
\begin{proof}
Let $i$ be a random block. Using the update \eqref{eq:update}:
\[
\begin{aligned}
\norm{w-\eta U_i\nabla_i f(w)-w^{\ast}}^{2}
&=
\norm{w-w^{\ast}}^{2}
-2\eta\ip{U_i\nabla_i f(w)}{w-w^{\ast}}
+\eta^{2}\norm{U_i\nabla_i f(w)}^{2}.
\end{aligned}
\tag{P1}
\]
Since $U_i^{\top}U_i=I_{n_i}$ and $U_iU_i^{\top}$ projects onto block $i$,
\[
\norm{U_i\nabla_i f(w)}^{2}= \norm{\nabla_i f(w)}^{2}.
\tag{P2}
\]
Take expectation over $i$.
Because each block is selected with probability $1/m$,
\[
\E[U_i\nabla_i f(w)]=\frac1m\sum_{i=1}^{m}U_iU_i^{\top}\nabla f(w)=\frac1m\nabla f(w).
\tag{P3}
\]
Similarly,
\[
\E\bigl[\norm{\nabla_i f(w)}^{2}\bigr]=\frac1m\sum_{i=1}^{m}\norm{\nabla_i f(w)}^{2}
\le \frac{L_{\max}}{m}\norm{\nabla f(w)}^{2},
\tag{P4}
\]
where the inequality follows from $\norm{\nabla_i f(w)}^{2}\le L_i\,\norm{\nabla_i f(w)}^{2}\le L_{\max}\norm{\nabla_i f(w)}^{2}$ and $\sum_i\norm{\nabla_i f(w)}^{2}=\norm{\nabla f(w)}^{2}$.

Plugging (P3)–(P4) into the expectation of (P1) yields
\[
\begin{aligned}
\E\bigl[\norm{w-\eta U_i\nabla_i f(w)-w^{\ast}}^{2}\bigr]
&=
\norm{w-w^{\ast}}^{2}
-2\eta\frac{1}{m}\ip{\nabla f(w)}{w-w^{\ast}}
+\eta^{2}\frac{L_{\max}}{m}\norm{\nabla f(w)}^{2}.
\end{aligned}
\tag{P5}
\]
By convexity (A1), $\ip{\nabla f(w)}{w-w^{\ast}}\ge f(w)-f(w^{\ast})$. Multiplying (P5) by $m$ and rearranging yields (L1). ∎
\end{proof}

\section*{Main Proof (Step‑by‑Step Derivation)}
We prove Theorem~\ref{thm:conv}.  
All expectations are taken with respect to the random block selections $\{i_t\}_{t=0}^{T-1}$.

\begin{enumerate}
\item[Step 1.] Apply Lemma~\ref{lem:dist} to the iterate $w_t$:
\[
\E_{i_t}\!\bigl[\norm{w_{t+1}-w^{\ast}}^{2}\mid w_t\bigr]
\le
\norm{w_t-w^{\ast}}^{2}
-2\eta\bigl(f(w_t)-f(w^{\ast})\bigr)
+\eta^{2}L_{\max}\norm{\nabla f(w_t)}^{2}.
\tag{S1}
\]

\item[Step 2.] Use the block‑coordinate smoothness \eqref{eq:smooth} with $d=-\eta\nabla_i f(w_t)$ and the inequality $ \norm{\nabla_i f(w_t)}^{2}\le L_i\bigl(f(w_t)-f(w^{\ast})\bigr)$ that follows from convexity and smoothness (standard descent lemma). Since $L_i\le L_{\max}$,
\[
\eta^{2}L_{\max}\norm{\nabla f(w_t)}^{2}
\le 2\eta\bigl(f(w_t)-f(w^{\ast})\bigr).
\tag{S2}
\]

\item[Step 3.] Substitute (S2) into (S1) to cancel the gradient‑norm term:
\[
\E_{i_t}\!\bigl[\norm{w_{t+1}-w^{\ast}}^{2}\mid w_t\bigr]
\le
\norm{w_t-w^{\ast}}^{2}
-2\eta\bigl(f(w_t)-f(w^{\ast})\bigr)
+2\eta\bigl(f(w_t)-f(w^{\ast})\bigr)
=
\norm{w_t-w^{\ast}}^{2}.
\tag{S3}
\]
Thus the expected squared distance is non‑increasing.

\item[Step 4.] Take total expectation and sum (S1) over $t=0,\dots ,T-1$:
\[
\begin{aligned}
\E\!\bigl[\norm{w_T-w^{\ast}}^{2}\bigr]
&\le
\norm{w_0-w^{\ast}}^{2}
-2\eta\sum_{t=0}^{T-1}\E\!\bigl[f(w_t)-f(w^{\ast})\bigr]
+\eta^{2}L_{\max}\sum_{t=0}^{T-1}\E\!\bigl[\norm{\nabla f(w_t)}^{2}\bigr].
\end{aligned}
\tag{S4}
\]

\item[Step 5.] Apply the bound (S2) again to the last term:
\[
\eta^{2}L_{\max}\norm{\nabla f(w_t)}^{2}
\le 2\eta\bigl(f(w_t)-f(w^{\ast})\bigr).
\tag{S5}
\]
Insert (S5) into (S4):
\[
\E\!\bigl[\norm{w_T-w^{\ast}}^{2}\bigr]
\le
\norm{w_0-w^{\ast}}^{2}
-2\eta\sum_{t=0}^{T-1}\E\!\bigl[f(w_t)-f(w^{\ast})\bigr]
+2\eta\sum_{t=0}^{T-1}\E\!\bigl[f(w_t)-f(w^{\ast})\bigr].
\tag{S6}
\]

\item[Step 6.] The two summations cancel, leaving
\[
\E\!\bigl[\norm{w_T-w^{\ast}}^{2}\bigr]\le \norm{w_0-w^{\ast}}^{2}\le R^{2}.
\tag{S7}
\]

\item[Step 7.] Return to (S1) and drop the non‑negative term $\eta^{2}L_{\max}\norm{\nabla f(w_t)}^{2}$:
\[
2\eta\bigl(f(w_t)-f(w^{\ast})\bigr)
\le
\norm{w_t-w^{\ast}}^{2}-\E\!\bigl[\norm{w_{t+1}-w^{\ast}}^{2}\mid w_t\bigr].
\tag{S8}
\]

\item[Step 8.] Take full expectation on both sides of (S8) and sum over $t=0,\dots ,T-1$:
\[
2\eta\sum_{t=0}^{T-1}\E\!\bigl[f(w_t)-f(w^{\ast})\bigr]
\le
\norm{w_0-w^{\ast}}^{2}
-\E\!\bigl[\norm{w_T-w^{\ast}}^{2}\bigr]
\le R^{2}.
\tag{S9}
\]

\item[Step 9.] Divide by $2\eta T$ and use convexity of $f$ together with Jensen’s inequality:
\[
\E\!\bigl[f(\bar w_T)-f(w^{\ast})\bigr]
\le
\frac{1}{T}\sum_{t=0}^{T-1}\E\!\bigl[f(w_t)-f(w^{\ast})\bigr]
\le
\frac{R^{2}}{2\eta T}.
\tag{S10}
\]

\item[Step 10.] Inequality (S10) is exactly the claimed bound (T1). ∎
\end{enumerate}

\section*{Rate Derivation (With Constants)}
From (T1) we have the explicit convergence rate
\[
\boxed{\;
\E\!\bigl[f(\bar w_T)-f(w^{\ast})\bigr]
\le
\frac{R^{2}}{2\eta\,T},
\qquad
0<\eta\le\frac{1}{L_{\max}}.
\;}
\]
If we choose the maximal admissible stepsize $\eta=1/L_{\max}$, the bound simplifies to
\[
\E\!\bigl[f(\bar w_T)-f(w^{\ast})\bigr]
\le
\frac{L_{\max}R^{2}}{2T}=O\!\left(\frac{1}{T}\right).
\]

\section*{Special Cases}
\begin{itemize}
\item \textbf{Uniform block sizes and $L_i=L$ for all $i$.} Then $L_{\max}=L$ and the rate reduces to the classic $O(1/T)$ for coordinate descent on smooth convex functions.
\item \textbf{Strongly convex $f$ with parameter $\mu>0$.} A standard modification of the above proof (adding the strong convexity inequality) yields a linear rate:
\[
\E\!\bigl[f(w_T)-f(w^{\ast})\bigr]\le
\left(1-\frac{\eta\mu}{m}\right)^{T}\bigl[f(w_0)-f(w^{\ast})\bigr].
\]
\item \textbf{Non‑uniform sampling.} If block $i$ is selected with probability $p_i>0$, the same analysis holds after replacing $L_{\max}$ by $\max_i L_i/p_i$ and $m$ by $\sum_i p_i^{-1}$.
\end{itemize}

\end{document}